% !TeX root = Tesis.tex
\chapter{Introducción}
\label{Ch:Intro}

Las ciencias de la computación evolucionan día a día. Su historia
a principios del siglo XX donde aparecen las primeras
computadoras,éstas eran máquinas que ocupaban gran espacio y no
contaban con entorno gráfico para programarlas todo se hacia de
manera de manual. Años más tarde aparecen las computadoras de
escritorio, posteriormente las laptop y en nuestros días los
dispositivos móviles. A partir de la aparición de los dispositivos
móviles la tecnología ha optado en construir entornos inteligentes
de tal manera que se pueda controlar el medio físico. De esta
manera nace la idea de crear entornos ubicuos.

\vspace*{0.4cm} Un entorno ubicuo es un área autónoma e
independiente, y que administra sus propios recursos y servicios,
intercambiando información con otras áreas para coordinarse entre
ellas y tomar decisiones en función del contexto actual del
usuario nómada \cite{Castro2010}. Un usuario nómada, es aquel que
se desplaza dentro del entorno ubicuo entre las diferentes áreas
tomando un rol diferente en cada una de ellas. Este le permite
acceder a la información y a los servicios que se ofrecen en estas
áreas.

\vspace*{0.4cm} A continuación se muestra la estructura general
que conforma un entorno ubicuo:\\

\vspace*{0.2cm}

\begin{figure}
  \begin{center}
    \includegraphics[width=7cm]{pictures//egentornoubicuo.png}\\
   \caption{Estructura general de un entorno ubicuo}\label{Fig1. Estructura general de un entorno ubicuo}
  \end{center}
\end{figure}

\newpage
En la imagen se muestra cuando un usuario nómada va a interactuar
en el entorno ubicuo y puede acceder a diferentes servicios
ofrecidos por estas áreas, como pueden ser: acceder a información
administrativa, personal, finanzas, recursos humanos; compartir
información entre los usuarios; acceder a Internet, hacer pagos,
ingresar a páginas web, hacer compras y videoconferencias con
otras entidades físicas relacionadas a la organización. En la
figura anterior se hace mención a estas áreas etiquetando como
área1, área2, área3 y una área en común denominada área de
administración de flujos de trabajo. A las áreas anteriores se
pueden añadir el control de los servicios físicos tales como;
impresión y escaneado de documentos, control de temperatura,
control de iluminación, control del sistema de seguridad, etc.\\

\vspace*{0.2cm}

La importancia de controlar el medio físico permite tener mejores
servicios optimizando el rendimiento de los mismos, por ejemplo:
solamente debe estar encendida una lámpara cuando exista personal
que requiera de dicho servicio, mientras tanto no es recomendable
que se tenga encendida la luz si no se tiene personal en dicha
área, he aquí la importancia de crear los entornos ubicuos, de tal
manera que permitan controlar el ambiente de manera óptima y
eficiente.\\

\vspace{0.2cm}

Dentro del entorno ubicuo el usuario nómada puede manipular los
servicios ofrecidos por cada área empleando dispositivos móviles,
a través del área de administración de flujos de trabajo, la cual
se encarga de controlar los servicios físicos.\\

\vspace{0.2cm} Pero ahora se preguntará ?`Cómo nacen los entornos
ubicuos?\\

\vspace{0.2cm}

La historia de los sistemas de cómputo nace a partir de 1970
cuando surgen los sistemas distribuidos. Estos son considerados
como una colección de computadoras independientes que aparecen
ante los usuarios del sistema, como una única
computadora\cite{Tanenbaum1996}. Los sistemas distribuidos forman
una LAN que permite a los usuarios acceder a la información de
manera remota. Con la aparición de las redes inalámbricas y las
computadoras portátiles,permiten acceder a la información de
manera móvil, se da una evolución de los sistemas distribuidos
hacia el cómputo móvil, que nace en la década de 1990. Los
sistemas de cómputo móvil permiten acceder de manera inalámbrica
desde cualquier dispositivo PDA, Laptop o cualquier teléfono
celular con capacidad de transmisión y recepción a múltiples
servicios como e-mail, bases de datos, sistemas ERP, redes
privadas, internet e intranet, en forma sencilla y rápida
\cite{computomovil}. Con el avance de la tecnología los
dispositivos portátiles tienden a ser más pequeños, incrementando
su mayor uso en aplicaciones, que no sólo permiten acceder a la
información sino a los servicios que se ofrecen en las redes
inalámbricas. La disminución en el tamaño de estos dispositivos
móviles hace que sean imperceptibles a la vista del usuario,
creando con ello los espacios inteligentes, dando origen así a la
idea de cómputo ubicuo. Mark Weiser visualiza esta idea de cómputo
ubicuo como: ``las tecnologías más profundas son las que
desaparecen... se tejen ellas mismas dentro de la vida cotidiana
hasta que se vuelven indistinguibles'', lo cual permite a los
usuarios acceder a la información en entornos denominados entornos
ubicuos.\\
\vspace{0.2cm}

Como ya he mencionado anteriormente los entornos ubicuos cuentan
con el área de administración de flujos de trabajo, el cual
trabaja de la siguiente manera:\\
\vspace{0.2cm}

La comunicación entre el dispositivo móvil y el área de
administración de flujos de trabajo se realiza mediante
comunicación bluetooth, mientras que la comunicación del área de
administración de flujos de trabajo y los diferentes sistemas de
control permite controlarlos mediante la comunicación USB que
todavía no existen para los entornos ubicuos pero son primordiales
para establecer la comunicación con los sistemas físicos que se
encuentren en el entorno. De aquí surge la necesidad de generar el
protocolo de comunicación vía USB para estos entornos ubicuos.\\

\vspace{0.2cm}

El crear estos entornos ubicuos debe contar con la ayuda de
diversas disciplinas de las ciencias abstractas. En este caso la
informática forma parte de las ciencias abstractas, y cuya
definición es el manejo automático de la información. En mi caso
como Informático me motivó desarrollar este tema de investigación
por que al crear un protocolo de comunicación se manejará
eficientemente los recursos computacionales. De igual forma
controlar de los equipos físicos de un entorno, resultado de ello,
que se utilicen de la mejor manera sin que se desperdicien, por el
mal uso que se le da.\\

\vspace{0.2cm}

El generar el protocolo de comunicación vía USB permitirá que un
entorno ubicuo pueda manipular los equipos físicos y automatizar
las áreas ofreciendo mejores servicios físicos a los usuarios
nómadas.\\

\vspace{0.2cm}

Un usuario nómada puede manipular los servicios ofrecidos por cada
área, empleando dispositivos móviles, a través del área de
administración de flujos de trabajo, misma que controla los
equipos físicos dentro del entorno a través del sistema de
control.\\

\vspace{0.2cm}

La importancia que juega el área de administración de flujos de
trabajo en un entorno ubicuo radica en que es el área en donde se
almacena la información para controlar los servicios físicos
(control de temperatura, ventilación, abrir y cerrar una ventana o
una puerta, etc.) y los servicios lógicos (aplicaciones web,
acceso a archivos, etc.) a los que puede tener acceso el usuario
nómada. Así que es importante contar el área de administración de
flujos de trabajo como un área común para todas las áreas de
manera que se cumpla con la definición del cómputo ubicuo, proveer
servicios a los usuarios en redes que se tejen en objetos de la
vida diaria, pero que sean imperceptibles a la vista del usuario.
Y la informática puede contribuir a este desarrollo automatizando
el flujo de información entre el área de administración de flujos
de trabajo y el sistema de control.\\

\vspace{0.2cm}

\newpage
\section{Planteamiento del Problema}
Los entornos hospitalarios ubicuos son áreas autónomas que
trabajan de forma independiente, pero tienen que ser controladas y
administradas por el área de administración de flujos de trabajo
(AAFT), que se encarga de ofrecer los servicios disponibles en la
red a los usuarios nómadas. Estos servicios pueden ser web o de
manipulación de servicios físicos. Para que el usuario nómada
pueda manipular los servicios físicos a través del área de
administración de flujos de trabajo, es necesario disponer de un
sistema de control, que se encargue de controlar dichos servicios.
El sistema de control es un conjunto de elementos que permiten
controlar un sistema (e.g. iluminación, ventilación, abrir y
cerrar puertas y ventanas, entre otros). Básicamente, un sistema
de control esta formado por sensores, actuadores, sistema de
acondicionamiento de señal y etapa de potencia.\\

\vspace{0.2cm}

Actualmente en los entornos hospitalarios ubicuos, no se
especifica el sistema de control, el área de administración de
flujos de trabajo, ni mucho menos la forma como se comunican ambos
sistemas. Se ha venido considerando que es suficiente tener un
área de administración de flujos de trabajo para manipular
directamente los servicios físicos, lo cual no es cierto, porque
se tiene la necesidad de tener un sistema de control entre el AAFT
y los sistemas físicos, el problema surge de la necesidad de crear
la comunicación entre el área de administrador de flujos de
trabajo y el sistema de control que es lo que no se tiene y es
primordial en los entornos hospitalarios ubicuos.\\

\vspace{0.2cm}

Como esta es la primera aportación a los entornos ubicuos todavía
no se tiene terminadas todas las partes que conforman los entornos
ubicuos, así es que la prueba para el protocolo de comunicación
vía USB consistirá en una computadora que simule al AFT y el
microcontrolador el sistema de control. Esto se realizará enviando
y recibiendo datos entre la computadora y el microcontrolador a
través de sus puertos USB mediante la entrada de datos por teclado
y visualización en pantalla de los resultados a través del monitor
un LCD respectivamente.\\

\vspace{0.2cm}

La importancia de crear los protocolos en el área de informática
radica en la facilidad que se obtendrá en el envió de la
información, lo cual será información que pasará por ciertas
reglas que estarán automatizadas para facilitar el acceso al área
de administrador de flujos de trabajo y al sistema de control,
mismo que la informática se encarga del manejo automático de la
información y con este aporte la informática se ve involucrado en
el manejo eficiente de los recursos computacionales con los que se
cuente para proveer los servicios que se requieran. Toda esta
investigación de los entornos hospitalarios ubicuos esta en pleno
desarrollo aún no se han concluido las demás etapas, por
consiguiente esta investigación quienes vallan a probar el
protocolo son investigadores que del área de cómputo ubicuo,
computo móvil y sistemas de control. Aunque es una investigación
que beneficiará al sector salud, en esta ocasión no se contará con
el apoyo del personal de esta área para probar este protocolo, por
que falta la interacción de otras áreas que se ven involucradas en
un entorno hospitalario ubicuo.\\

\vspace{0.2cm}

\newpage
\section{Objetivo General}

Desarrollar un protocolo de comunicación USB que posibilite la
comunicación entre el administrador de flujos de trabajo y el
sistema de control en un entorno hospitalario ubicuo.\\

\vspace{0.2cm}

\section{Objetivos Específicos}

\begin{itemize}
    \item Extraer los requisitos de la manera en que se almacenan los datos en el administrador de flujos de trabajo.
    \item Entregar la documentación de la forma en que trabaja el área de control de la computadora.
    \item Obtener la configuración correcta del puerto USB en el administrador de flujos de trabajo para enviar/recibir datos.
    \item Desarrollar la interfaz grafica de configuración del puerto USB con las librerías de Windows archivos (*.dll) y la API de Java para facilitar el envió/recepción de los datos.
    \item Entregar el manual de configuración del protocolo USB para configurar correctamente el administrador de flujos de trabajo de manera que se pueda implementar la comunicación bidireccional del puerto USB al área de control.
\end{itemize}

\section{Supuestos}

A través del protocolo de comunicación USB desarrollado se
comunican correctamente el área de administración de flujos de
trabajo y el sistema de control de un entorno hospitalario
ubicuo.\\

\vspace{0.1cm}

El protocolo de comunicación vía USB manipulara de mejor manera el
área de control de manera que no necesite tanta configuración para
el envió/recepción de datos.\\

\vspace{0.1cm}

El protocolo de comunicación vía USB permite tener el control de
los datos, para manipular eficazmente los dispositivos físicos a
través del sistema de control.\\

\vspace{0.1cm}

Los datos almacenados en el administrador de flujos de trabajo
permiten tener la correcta administración de los equipos físicos a
controlar y de los servicios lógicos que se pretenden ofrecer al
usuario.\\

\vspace{0.2cm}

\section{Hipótesis}

El protocolo de comunicación USB permite una comunicación de alta
velocidad que permite ser medida igual a la velocidad del
procesador para la transferencia de la información a través del
puerto USB 2.0 con el área de control.\\

\vspace{0.2cm}

\newpage
\section{Justificación}

En la actualidad, se cuenta con entornos hospitalarios donde la
obtención de información de los distintos servicios es
relativamente lenta, por lo que no se puede realizar una buena
administración física de las instalaciones del hospital.\\

\vspace{0.2cm}

En un hospital se tienen diversos usuarios, entre ellos están los
doctores, enfermeras, pacientes, personal administrativo, etc.
cada uno de ellos necesita acceder a diversos servicios físicos y
lógicos dependiendo del lugar donde se encuentre. Para ello se
debe contar con un sistema de control que visualice el entorno
físico y envíe datos al área de administración de flujos de
trabajo de un entorno hospitalario ubicuo. Este flujo de
información debe ser bidireccional, para que el usuario, a través
del área de administración de flujos de trabajo, pueda manipular
los diferentes servicios, y este es precisamente el problema que
se quiere resolver con el desarrollo de este proyecto.\\

\vspace{0.2cm}

El beneficio que aportará este proyecto de tesis es empezar a
construir un entorno hospitalario ubicuo, ya que cada día en los
hospitales se esta empezando a automatizar áreas y servicios, y se
esta invirtiendo en infraestructura y en tecnología para equiparlo
de la mejor calidad, mismo que al crear un entorno hospitalario
ubicuo da inicio a la creación de entornos hospitalarios
ubicuos.\\

\vspace{0.2cm}

Por ejemplo, en el distrito de Miahuatlán de Porfirio Díaz Oaxaca
es una región donde se tiene un creciente número de habitantes y
se cuenta con poco personal médico para atender a los pacientes,
por lo tanto deben de automatizarse algunas áreas de hospital, la
solución ideal es empezar a crear entornos hospitalarios ubicuos,
por ejemplo automatizar el área de administración del hospital. El
beneficio que se obtendrá será un mejor servicio de salud a la
población.Algunos ejemplos:\\

\begin{itemize}
    \item Agilizar el tiempo de consulta para la adqusición del
    expediente clinico del paciente.
    \item Ahorro de energía eléctrica en la infraestructura del
    hospital, cuando no se ocupe dicha área.
    \item Cuidar del medio ambiente al no utilizar energía cuando
    no es necesario.
\end{itemize}

\vspace{0.2cm}

El crear un entorno hospitalario ubicuo beneficiará a médicos,
enfermeras, pacientes, personal administrativo, etc. ya que
facilitará el manejo de la información y el acceso a ello será más
ágil, este proyecto es la primera aportación que se dará al crear
un entorno hospitalario ubicuo, ya que en un futuro estos entornos
mejoraran la eficiencia y eficacia para proporcionar mejores
oportunidades de acceso a la salud pública equipando con la mejor
y última tecnología para los hospitales del país.\\

\vspace{0.2cm}

Actualmente, se están haciendo inversiones para equipar hospitales
con las mejores tecnologías de la información, a fin de mejorar la
calidad de los servicios que se ofrecen. El implementar un entorno
hospitalario ubicuo permite automatizar las distintas áreas de un
hospital logrando con ello que se trabaje de mejor manera y
ofreciendo un servicio de calidad a los usuarios, donde juega un
papel importante la comunicación entre el área de
administración de flujos de trabajo y el sistema de control.\\

\vspace{0.2cm}

Esta investigación es importante por que forma parte del
desarrollo de los entornos hospitalarios ubicuos. Establecer un
protocolo de comunicación entre el administrador de flujos de
trabajo y el sistema de control permitirá tener toda la
automatización antes mencionada, y con la información que se
almacene en el administrador de flujos de trabajo para controlar
todas las áreas, pero sobre todo radica en la correcta
comunicación que se debe implementar, y en esto se apoya de la
Informática, del manejo
automático de la información.\\

\vspace{0.2cm}

\section{Metas}
\begin{itemize}
    \item Obtener el software del protocolo de comunicación USB para el área de administración de flujos de trabajo como para el sistema de control.
    \item Documentar las características, reglas y demás especificaciones del protocolo de comunicación USB.
    \item Realizar un manual del programador que explique la configuración del puerto USB.
    \item Realizar un manual de usuario donde se explique la instalación y configuración del software del protocolo de comunicación.
    \item Diseñar los diagramas al flujo de información entre el AAFT y el sistema de control.
    \item Validar el protocolo de comunicación USB conectando al puerto USB a un cable USB y éste a un LCD y que simule la transferencia de datos enviando datos desde el administrador de flujos de trabajo y el sistema de control.
\end{itemize}

\newpage
\section{Marco teórico}

\subsection{Cómputo ubicuo}
En la actualidad los usuarios de las tecnologías de la información
y comunicación están pasando de la era de la computadora personal
PC (Personal Computer), a la era de los dispositivos móviles (e.g.
Smartphones, computadoras portátiles y PDAs), donde por lo regular
se tiene acceso a diferentes servicios y éstos se proporcionan de
manera inalámbrica.\\

\vspace{0.2cm}

Estos dispositivos móviles cada día tienden a hacerse más pequeños
y con mayor capacidad de procesamiento de información,
ofreciéndose mayores prestaciones y servicios a través de las
redes inalámbricas.\\

\vspace{0.2cm}

La idea de Mark Weiser, considerado como el padre del cómputo
ubicuo, consiste en la desaparición física de la computadora y en
la incorporación de capacidades de conmutación y comunicación en
la mayoría de objetos físicos cotidianos, con el fin de formar una
gran red de dispositivos interconectados, cuyo objetivo es crear
entornos físicos inteligentes mediante el uso de sistemas
computacionales, que estén disponibles pero que se vuelvan
invisibles a los usuarios\cite{Castro2010}. La idea anterior
define lo que se
conoce como cómputo ubicuo.\\

\vspace{0.2cm}

Un espacio inteligente para que funcione correctamente, debe estar
orientado hacia lo que definió Weiser como cómputo ubicuo,
surgiendo así los entornos ubicuos que constan de muchas
computadoras colocadas en un espacio físico, unidas por una red de
alta velocidad y capaces de detectar personas, objetos y
dispositivos.\\

\vspace{0.2cm}

\subsection{Áreas Autónomas}

Un entorno ubicuo está constituido por áreas autónomas e
independientes, se llaman de esta manera por que administran sus
propios recursos y servicios intercambiando información con otras
áreas para coordinarse entre ellas y toma decisiones en función
del contexto actual del usuario\cite{Tanenbaum1996}.\\

\vspace{0.2cm}

Dentro de las áreas autónomas, el usuario tiende a volverse
nómada. El usuario puede manipular distintas áreas y no está
limitado a una de ellas. Así nace el concepto de usuario nómada.\\

\subsection{Usuario nómada}

El usuario nómada es aquel que se desplaza dentro de las áreas
autónomas del entorno ubicuo tomando un rol distinto en cada una
de ellas, lo que le permite acceder a la información y a los
servicios que se ofrecen en dichas áreas.\\

\vspace{0.2cm}

\subsection{Área de administración de flujos de trabajo (AAFT)}

Como he mencionado anteriormente, los entornos ubicuos están
constituidos por áreas autónomas independientes de las otras
áreas, pero para que puedan coordinarse e intercambiar
información, existe un área común llamada área de administración
de flujos de trabajo.\\

\vspace{0.2cm}

El área de administración de flujos de trabajo está formada por
los sistemas de flujos de trabajo (workflows) enfocados a la
definición y construcción de los procesos de las organizaciones.
Estas se determinan en términos de tareas\cite{articulodeZur} .
Los usuarios son definidos como actores, los que llevan a cabo un
conjunto de tareas, siguiendo un orden parcial, que define el
momento en que
deben ser analizadas.\\

\vspace{0.2cm}

Entonces, para que exista comunicación en un entorno ubicuo debe
existir el área de administración de flujos de trabajo, que se
encarga de coordinar la comunicación entre las diversas áreas del
entorno.\\

\vspace{0.2cm}

\subsubsection{Sistemas de control}

Para poder acceder a los servicios físicos que se ofrecen en las
diversas áreas autónomas, es necesario disponer de uno o varios
sistemas de control que permitan manipular los servicios físicos.
Un elemento físico a ser controlado se denomina``Planta'', por
ejemplo: la temperatura, la iluminación de una sala, puertas,
ventanas, cafeteras eléctricas, etc.\\

\vspace{0.2cm}

Una combinación de componentes o dispositivos que actúan juntos
para alcanzar un cierto objetivo se denomina ``Sistema''. Un
sistema de control es un sistema que consta de muchos componentes
como: sensores, actuadores, plantas, etc. El objetivo del sistema
de control es llevar al sistema de un estado inicial a un estado
final, en un tiempo finito\cite{libroScontrol}.\\

\vspace{0.2cm}

\subsection{Puertos USB}
En este proyecto el AAFT y el sistema de control, se van a
conectar entre sí utilizando puertos USB.\\

\vspace{0.2cm}

USB es el acrónimo de Universal Serial Bus \cite{librodeUSB}. El
USB fue diseñado para establecer una interfaz de comunicación con
muchos
tipos de periféricos de las interfaces.\\

\vspace{0.2cm}

El puerto USB nace a partir de 1996 con el USB 1.0, fecha en la
que es creado por un conjunto de empresas de las cuales se
destacan IBM, Intel, Northern Telecom, Compaq, Microsoft, Digital
Equipment Corporation y NEC\cite{puertosES}. Dos años más tarde,
en 1998, se lanza el USB 1.1, que agrega un Interruptor de salida.
Y posteriormente, en abril del año 2000, se lanza al mercado el
USB
2.0 que agrega opciones de uso de alta velocidad.\cite{librodeUSB}\\

\vspace{0.2cm}

Cuando se inició el diseño del USB se tenía en mente eliminar la
necesidad de adquirir tarjetas separadas para poner en los puertos
ISA y PCI y mejorar las capacidades plug and play, permitiendo a
estos dispositivos ser conectados y desconectados al sistema sin
la necesidad de reiniciar el sistema operativo.\\

\vspace{0.2cm}

El auge que tuvo el USB en sus primeros años de creación, fue
gracias al incremento en la velocidad de transmisión y adquisión
de datos del sistema operativo hacia el USB, lo que llevó a
reemplazar los teclados y mouse con puertos PS/2, llegando hasta
el punto de que las placas base modernas carecen ya de dicho
puerto y ahora solamente cuentan con puertos USB.\\

\vspace{0.2cm}

El USB es una de las mejores alternativas para el envío,
transmisión y recepción de información de un dispositivo hardware
al sistema operativo. Para ello existen cuatro clasificaciones de
la velocidad de transferencia de los
dispositivos USB:\\

\begin{itemize}
    \item Baja velocidad: Utilizado por el USB 1.0, con una tasa de transferencia de 1.5 Mbps o 192 KB/s.
    \item Velocidad completa: Se implementó en el USB 1.1 con una tasa de transferencia de hasta 12 Mbps o 1.5 MB/s.
    \item Alta velocidad: Actualmente se encuentra en esta categoría el USB 2.0, con una tasa de transferencia de hasta 480 Mbps o 60 MB/s.
    \item Super alta velocidad: Esta velocidad se implementará en el USB 3.0 con una tasa de transferencia de datos de hasta 4.8 Gbps o 600
    MB/s.\cite{librodeUSB}
\end{itemize}

\vspace{0.2cm}

\subsection{Microcontrolador}

Un microcontrolador es un microprocesador con memoria de datos y
de programas, puertos de entrada y salida, interruptores externos
e internos, todo integrado en un solo chip.\cite{libroPIC}\\

\vspace{0.2cm}

Dentro de las aplicaciones del microcontrolador se encuentran: PC,
impresoras láser, máquinas de fax y teléfono inteligentes. También
son utilizados en DVDs, video juegos y máquinas de lavado.\\

\vspace{0.2cm}
